% Folder 144, batch 08 — archivists' pages 141–160.
% Pass: Opus 5.5 (claude-opus-5-5), 2026-10-03 — first pass, unchecked against the pages by a human.
%
% Skipped as unrelated versos (administrative documents, computer listings):
% 150, 151, 152, 154, 156, 158, 160.
% The pencilled archivists' numbers on these sheets read one higher than the
% batch numbering used here (page 141 of the batch bears "142"); \page{} follows
% the batch/PDF numbering set up for this folder.
\documentclass[11pt,a4paper]{article}
\input{../preamble/grothendieck}

\folder{144}
\batch{8}
\pages{141}{160}
\foldertitle{[Suite autour de Teichmüller dont] Teichmülleries, 1981-1982 : notes manuscrites (1981-1983, s.d.), lettres (1981, s.d.).}
\dating{1981-1983}
\watermark{Édition de démonstration}

\begin{document}

\page{141}
\section{1) Groupe de Teichmüller spécial étendu du cylindre}
\note{titre de sa main, en tête de la page, numéroté « 1) ».}

Considérons le cylindre standard
\[
(1)\qquad X = \mathbb{U}\times\mathbb{I}, \quad\text{où}\quad
\mathbb{U}=\{z\in\mathbb{C} \mid |z|=1\},\quad \mathbb{I}=[-1,+1],
\]
donc
\[
(2)\qquad \partial X = \mathbb{U}\times\partial\mathbb{I}
= \mathbb{U}\times\{-1,+1\} = \mathbb{U}_{-1}\sqcup\mathbb{U}_{1}.
\]
Choisissons le revêtement universel standard de $\mathbb{U}_\varepsilon$ ($\varepsilon\in\{\pm1\}$), par
\[
(3)\qquad
\begin{cases}
\widetilde{\mathbb{U}}_\varepsilon \overset{\text{déf}}{=} \mathbb{R}_\varepsilon \longrightarrow \mathbb{U}_\varepsilon\\
(t,\varepsilon)\longmapsto (\exp(2i\pi t),\varepsilon).
\end{cases}
\]
Considérons le groupe « dihédral torique » (\uncertain{$\simeq O(2,\mathbb{R})$})
\[
(4)\qquad D_{\mathbb{U}} \overset{\text{déf}}{=} \{1,\sigma\}.\mathbb{U}\qquad(\tfrac12\text{ direct})
\]
\note{la parenthèse « ($\simeq$ O(2,R)) » est ajoutée au-dessus de la ligne et reliée à « diédral torique » ; devant le O une lettre surchargée, peut-être un S biffé.}
où $\sigma$ est la conjugaison complexe de $\mathbb{C}$, opérant sur $\mathbb{U}$ par $z\mapsto z^{-1}$, nous choisissons comme revêtement universel de ce groupe le groupe « dihédral réel »
\[
(5)\qquad \widetilde{D}_{\mathbb{U}} \overset{\text{déf}}{=} D_{\mathbb{R}} = \{1,\sigma\}.\mathbb{R},
\]
\uncertain{avec} l'homomorphisme de revêtement
\[
(6)\qquad \sigma^\nu.t \longmapsto \sigma^\nu \exp(2i\pi t),
\]
\marginal{$\nu\in\mathbb{Z}/2\mathbb{Z}$}
de sorte qu'on a une suite exacte
\[
(7)\qquad 1\longrightarrow\mathbb{Z}\longrightarrow D_{\mathbb{R}}\longrightarrow D_{\mathbb{U}}\longrightarrow 1.
\]
Nous faisons opérer $D_{\mathbb{R}}$ sur $\mathbb{R}_\varepsilon=\widetilde{\mathbb{U}}_\varepsilon$ par

\page{142}
\[
(8)\qquad (\sigma^\nu.u).t_\varepsilon = (-1)^\nu(t+u)_\varepsilon
\]
(N.B. $t_\varepsilon=(t,\varepsilon)$, $\varepsilon\in\{\pm1\}$).
Cette opération passe au quotient en une opération de $D_{\mathbb{U}}$ sur \struck{\ill{}} $\mathbb{U}_\varepsilon$, et cette dernière est donnée par
\[
(9)\qquad (\sigma^\nu.\theta).z_\varepsilon = (z\theta)^{(-1)^\nu} =
\begin{cases}
z\theta & \text{si } \nu\equiv0\ (2)\\
z^{-1}\theta^{-1} & \text{si } \nu\equiv1\ (2).
\end{cases}
\]
Ainsi le groupe $D_{\mathbb{R}}\times D_{\mathbb{R}}$ opère sur
\[
\widetilde{\partial X} \overset{\text{déf}}{=} \widetilde{\mathbb{U}}_{-1}\sqcup\widetilde{\mathbb{U}}_{1} = \mathbb{R}_{-1}\sqcup\mathbb{R}_{1}
\]
via ses actions sur les deux sommandes, et cette action \ill{} passe aux quotients en une action de $D_{\mathbb{U}}\times D_{\mathbb{U}}$ sur \struck{\ill{}} $\partial X=\mathbb{U}_{-1}\sqcup\mathbb{U}_{1}$. Nous allons prolonger cette action de $D_{\mathbb{R}}\times D_{\mathbb{R}}$ sur $\widetilde{\partial X}$ en une action sur
\[
(10)\qquad \widetilde{X} \overset{\text{déf}}{=} \mathbb{R}\times\mathbb{I}
\]
qui revêt $X$ par
\[
(11)\qquad (t,s)\longmapsto(\exp 2i\pi t, s)\ ;
\]
\struck{Pour ceci} \add{plus commodément}, c'est le sous-groupe
\[
(12)\qquad
\begin{gathered}
D_{\mathbb{R}\times\mathbb{R}} \overset{\text{déf}}{=} \{1,\sigma\}.(\mathbb{R}\times\mathbb{R}) \hookrightarrow D_{\mathbb{R}}\times D_{\mathbb{R}}\qquad(\tfrac12\text{ direct})\\
\sigma^\nu.(u_{-1},u_1)\longmapsto(\sigma^\nu.u_{-1},\sigma^\nu.u_{+1})
\end{gathered}
\]
qu'on va faire opérer, de façon à ce que ces opérations passent au quotient en une opération de $D_{\mathbb{R}\times\mathbb{R}}$ sur $X$ lui-même. On

\page{143}
désignons, pour un élément
$(u_{-1},u_1)\in\mathbb{R}\times\mathbb{R}$,
\note{ici et plus bas l'indice de $D$ porte sur le premier $\mathbb{R}$ une marque qu'on peut prendre pour un astérisque ; on lit $\mathbb{R}\times\mathbb{R}$, conformément à (12).}
par
\[
(13)\qquad s\longmapsto u_s
\]
l'application affine de $\mathbb{R}$ dans $\mathbb{R}$ \struck{déf} qui prend la valeur $u_{-1}$ pour $s=-1$, la valeur $u_1$ pour $s=+1$, \struck{\ill{}} donc
\[
(14)\qquad u_s = \frac{u_1-u_{-1}}{2}\,s + \frac{u_1+u_{-1}}{2},
\]
et on pose
\[
(15)\qquad
\begin{cases}
(u_{-1},u_1).(t,s) \overset{\text{déf}}{=} (t+u_s,s)\\
\sigma.(t,s) \overset{\text{déf}}{=} (-t,s)
\end{cases}
\]
d'où une opération de $D_{\mathbb{R}\times\mathbb{R}}$ sur $\widetilde{X}$, compatible d'ailleurs avec la projection
\[
\widetilde{X}=\mathbb{R}\times\mathbb{I}\longrightarrow\mathbb{I},\qquad (t,s)\longmapsto s
\]
et l'opération triviale sur $\mathbb{I}$. L'opération passe bien au quotient en une opération de $D_{\mathbb{R}\times\mathbb{R}}$ sur $X$ lui-même, compatible \uncertain{avec} la projection
\[
X=\mathbb{U}\times\mathbb{I}\longrightarrow\mathbb{I},\qquad (z,s)\longmapsto s
\]
et l'opération triviale sur $\mathbb{I}$.

Considérons le sous-groupe de $D_{\mathbb{R}\times\mathbb{R}}$ qui opère trivialement sur $\partial X=\mathbb{U}\times\{-1,+1\}$, c'est le sous-groupe
\[
\mathbb{Z}\times\mathbb{Z}\subset D^{\circ}_{\mathbb{R}\times\mathbb{R}}.
\]

\page{144}
Considérons deux éléments
\[
(16)\qquad \sigma^\nu u_{\cdot}=\sigma^\nu.(u_{-1},u_1),\qquad \sigma^{\nu'}.u'_{\cdot}=\sigma^{\nu'}.(u'_{-1},u'_1)
\]
de $D_{\mathbb{R}\times\mathbb{R}}$, on voit de suite qu'ils ont même effet sur $\widetilde{X}$ \struck{\ill{}} (et \uncertain{même} seulement sur $\partial\widetilde{X}=\widetilde{\partial X}$) ssi ils sont égaux : $D_{\mathbb{R}\times\mathbb{R}}$ \emph{opère fidèlement} sur $\widetilde{X}$, et \uncertain{même} sur $\partial\widetilde{X}=\widetilde{\partial X}$.

Pour que les éléments (16) aient même action sur $X$, il faut et il suffit qu'on ait
\[
(17)\qquad \nu=\nu' \quad\text{et}\quad u_s-u'_s\in\mathbb{Z}\quad \forall s\in\mathbb{I}
\]
(la première condition \add{$\nu=\nu'$} exprime d'ailleurs que les deux opérations, envisagées \ill{}, ont même effet sur l'orientation de $X$). La condition $u_s-u'_s\in\mathbb{Z}$ $\forall s\in\mathbb{I}$ s'explique par
\[
(18)\qquad
\begin{cases}
\dfrac{u_1-u_{-1}}{2}=\dfrac{u'_1-u'_{-1}}{2} \quad\text{i.e.}\quad u_1-u_{-1}=u'_1-u'_{-1}\\[2ex]
\dfrac{u_1+u_{+1}}{2}\equiv\dfrac{u'_1+u'_{-1}}{2}\ (\mathbb{Z}) \quad \text{\struck{\ill{}}}
\end{cases}
\]
\note{au numérateur de gauche de la seconde ligne, le manuscrit porte $u_{+1}$ ; on attend $u_{-1}$.}

\struck{Écrivons} \add{Prenons} comme nouveaux paramètres, pour repérer un élément $\sigma^\nu(u_{-1},u_{+1})$ de $D_{\mathbb{R}\times\mathbb{R}}$, au lieu de ($\nu$ et) $u_{-1}$ et $u_{+1}$, mais les coefficients
$a=\dfrac{u_1-u_{-1}}{2}$ et $b=\dfrac{u_1+u_{-1}}{2}$ de (14), on pose
\[
(19)\qquad f_{a,b}\overset{\text{déf}}{=}(u_{-1},u_1)\quad\text{où}\quad
\begin{cases} u_{-1}=-a+b\\ u_1=a+b.\end{cases}
\]
Alors \ill{}
\[
(20)\qquad f_{a,b}=f_{a',b'} \iff a=a',\ b\equiv b'\ (\mathbb{Z})
\]
\note{il faut comprendre : même action sur $X$. Devant (20) un signe biffé.}

\page{145}
Ainsi, le groupe quotient de $D_{\mathbb{R}\times\mathbb{R}}$ qui opère fidèlement sur $X$ est le groupe
\[
(21)\qquad D_{\mathbb{R}\times\mathbb{U}} \overset{\text{déf}}{=} \{1,\sigma\}.(\mathbb{R}\times\mathbb{U})\qquad(\tfrac12\text{ direct})
\]
où $\sigma$ opère \add{encore} sur $\mathbb{R}\times\mathbb{U}$ par symétrie
\[
\sigma(u,\theta)=(-u,\theta^{-1}).
\]
On fait opérer $D_{\mathbb{R}\times\mathbb{U}}$ sur $X$ par
\[
(22)\qquad
\begin{cases}
\sigma(z,s)=(z^{-1},s)\\
g_{a,\theta}(z,s)\overset{\text{déf}}{=}(z\exp(2i\pi as)\,\theta,\ s)\qquad (a\in\mathbb{R},\ \theta\in\mathbb{U})
\end{cases}
\]
et cette action se relève en une action de $D_{\mathbb{R}\times\mathbb{R}}$ sur $\widetilde{X}$ par
\[
(23)\qquad
\begin{cases}
\sigma(t,s)\overset{\text{déf}}{=}(-t,s)\\
f_{a,b}(t,s)\overset{\text{déf}}{=}(t+(as+b),\ s)\qquad (a\in\mathbb{R},\ b\in\mathbb{R}).
\end{cases}
\]
\struck{Considérons} \add{On fera opérer} le groupe
\[
(24)\qquad \{\pm1\}\times D_{\mathbb{R}\times\mathbb{U}}
\]
\marginal{($\tfrac12$ direct, via op.\ de $\pm1$ sur $D_{\mathbb{R}\times\mathbb{U}}$ par chgt de signe sur le facteur $\mathbb{R}$)}
sur $X$, par
\[
(25)\qquad (\varepsilon\times\sigma^\nu.g_{a,\theta})(z,s)\overset{\text{déf}}{=}\bigl((z\exp(2i\pi as)\,\theta)^{(-1)^\nu},\ \varepsilon s\bigr) ;
\]
\note{après $z$, un exposant raturé.}
cette opération se réalise en une op.\ de
\[
(26)\qquad \{\pm1\}\times D_{\mathbb{R}\times\mathbb{R}}
\]
\marginal{($\tfrac12$ direct via op.\ de $\pm1$ sur $D_{\mathbb{R}\times\mathbb{R}}$ par chgt de signe sur le facteur $a$ de $f_{a,b}$ …)}
sur $\widetilde{X}$, par
\[
(27)\qquad (\varepsilon\times\sigma^\nu f_{a,b})(t,s)\overset{\text{déf}}{=}\bigl((-1)^\nu(t+(as+b)),\ \varepsilon s\bigr).
\]

\page{146}
Considérons le sous-groupe de \struck{\ill{}} $D_{\mathbb{R}\times\mathbb{U}}$ formé des $g_{a,\theta}$ qui \struck{ont} sur $(\partial X)_{-1}$ et $(\partial X)_{+1}$ sont des \struck{rotations} op.\ d'ordre fini, ce qui s'exprime par les conditions
\[
\theta\exp(-2i\pi a),\ \theta\exp 2i\pi a\in\mu_\infty(\mathbb{C})
\]
\marginal{(ens.\ des él.\ d'ordre fini de $\mathbb{U}$)}
soit encore
\[
(28)\qquad \theta\in\mu_\infty(\mathbb{C}),\quad a\in\mathbb{Q}.
\]
Ce sous-groupe est stable par l'opération de $\{\pm1\}\times\{1,\sigma\}$, opérant de la façon explicitée plus haut, \ill{} \ill{} un sous-groupe du groupe de (24), formé \add{justement} par les opérations \add{(\uncertain{identité que})} qui sur $\partial X$ sont d'ordre fini. Ce sous-groupe \struck{\ill{}} s'identifie à :
\[
(29)\qquad G=(\mu\times\mu).(\mathbb{Q}\times\mu_\infty)\qquad \tfrac12\text{ direct}
\]
où
\[
\mu=\{\pm1\},\qquad \mu_\infty=\mu_\infty(\mathbb{C})\simeq\mathbb{Q}/\mathbb{Z}\qquad
(\exp 2i\pi q \longleftarrow q \bmod \mathbb{Z}),
\]
et où $\mu\times\mu$ opère sur $\mathbb{Q}\times\mu_\infty$ par symétries sur les deux facteurs. On identifie bien $\sigma$ à $(-1,-1)\in\mu\times\mu$. On pose
\[
(30)\qquad \widetilde{G}=(\mu\times\mu).(\mathbb{Q}\times\mathbb{Q})\qquad\tfrac12\text{ direct}
\]
d'où la suite exacte
\[
(31)\qquad
\begin{gathered}
1\longrightarrow\mathbb{Z}\longrightarrow\widetilde{G}\longrightarrow G\longrightarrow 1\\
n\longmapsto f_{0,n},\qquad (\varepsilon,\varepsilon').f_{a,b}\longmapsto(\varepsilon,\varepsilon').g_{a,\exp 2i\pi b}
\end{gathered}
\]
\note{dans (29), après « $G=$ », quelques signes biffés, illisibles ; devant $\widetilde{G}$, dans (30) et (31), un signe biffé. Dans $n\mapsto f_{0,n}$ le premier indice est mal formé (on lirait aussi $1$) ; $0$ est ce que demande l'argument.}

\page{147}
L'opération de $\widetilde{G}$ sur $\widetilde{X}$ est donnée par
\[
(32)\qquad \bigl((\varepsilon,\varepsilon').f_{a,b}\bigr)(t,s)=\bigl(\varepsilon'(t+(as+b)),\ \varepsilon\varepsilon's\bigr)=\varepsilon'\bigl(t+(as+b),\ \varepsilon s\bigr)
\]
\[
(\varepsilon,\varepsilon'\in\mu,\quad a,b\in\mathbb{Q},\quad t\in\mathbb{R},\quad s\in\mathbb{I}),
\]
\note{dans la marge gauche, devant (32), une ligne biffée, illisible.}
elle passe au quotient en l'action de $G$ sur $X$
\[
(33)\qquad \bigl((\varepsilon,\varepsilon').g_{a,\theta}\bigr)(z,s)=\bigl((z\theta\exp 2i\pi as)^{\varepsilon'},\ \varepsilon\varepsilon's\bigr)
\qquad(\varepsilon,\varepsilon'\in\mu,\ a\in\mathbb{Q},\ \theta\in\mu_\infty).
\]
Pour l'opération $(\varepsilon,\varepsilon').g_{a,b}$, $\varepsilon$ est le caractère d'orientation (l'orientation conservée ou renversée suivant que $\varepsilon=+1$ ou $\varepsilon=-1$), $\varepsilon\varepsilon'$ est le signe qui indique l'action sur $\pi_0(\partial X)$ (« ensemble des trous »). L'action de $G$ sur $\partial X$ s'obtient en \uncertain{oubliant}, le signe $\varepsilon\varepsilon'$ pris \struck{et par les} indiquant s'il y a transposition entre les deux comp.\ connexes, et \struck{par les} le signe $\varepsilon'$ pris indiquant s'il faut composer \uncertain{avec} une conjugaison complexe, par les deux rotations
\[
\theta\exp(-2i\pi a),\quad \theta\exp(2i\pi a).
\]
\note{la suite de la page, jusqu'en haut de la page suivante, est barrée de traits obliques.}
\struck{Exprimons que l'\ill{}}
\[
\text{\struck{$(34)\quad \theta\exp(-2i\pi a)=-1,\quad \theta\exp 2i\pi a=-1$}}
\]
\struck{on trouve les deux} \add{\struck{types de}} \struck{solutions}
\[
\text{\struck{$(35)\quad \theta=1,\ a\in\tfrac12+\mathbb{Z}\quad;\quad \theta=-1,\ a\in\mathbb{Z}$}}
\]

\page{148}
\struck{ou plus simplement, si $\theta=\exp 2i\pi b$, en termes de $a$ et $b$}
\[
\text{\struck{$(36)\quad a\equiv\tfrac12,\ b\equiv1\ (\mathbb{Z})\quad\text{ou}\quad a\equiv1,\ b\equiv\tfrac12\ (\mathbb{Z}).$}}
\]
\struck{Nous nous intéresserons plus particulièrement aux sous-groupes}
\note{fin du passage barré.}

Nous nous intéresserons au sous-groupe $G_0$ de $G$ correspondant à la valeur $\theta=1$ ; on écrira $g_a$ au lieu de $g_{a,\theta}$, donc
\[
(34)\qquad g_a(z,s)=(z\exp 2i\pi as,\ s).
\]
\note{les numéros (34) et (35) sont repris après la rature.}
Pour que cette opération induise l'identité sur $\partial X$, il faut et il suffit qu'on ait
\[
(35)\qquad \exp 2i\pi a=1\quad\text{i.e.}\quad a\in\mathbb{Z}.
\]
Nous nous intéresserons au sous-groupe
\note{la phrase s'arrête ici ; le reste de la page est blanc.}

\page{149}
\note{changement d'encre (brune) : une nouvelle série de notes commence ici, sans lien explicite avec la numérotation (1)–(36) qui précède.}

\emph{Cylindre} \struck{\ill{}} $X$

trois ensembles à 2 éléments
a) $\pi_0\,\partial X$ — $\varepsilon_0(X)$

b) ens.\ des deux générateurs de $\pi_1(X)\simeq H_1(X)$ — $\varepsilon_1(X)$

c) une des deux orientations de $X$ — $\omega(X)$
\note{les trois noms $\varepsilon_0(X)$, $\varepsilon_1(X)$, $\omega(X)$ sont écrits en colonne à droite, séparés de la liste par un trait vertical.}
On a une iso.\ canonique
\[
\omega(X)\simeq\varepsilon_0(X)\wedge\varepsilon_1(X).
\]
\drawing{à gauche, un cylindre vertical ; sur le cercle du bas, une flèche tangente et une flèche verticale montant vers l'intérieur, marquant une orientation.}
Si on a une compos.\ $S$ du bord ($\alpha\in\varepsilon_0(X)$) et une orientation du bord (\struck{\ill{}} $\varepsilon_1(X)\simeq$ un des deux génér.\ de $H_1(S)$), on a une orientation de $X$ : celle qui induit $\beta$ sur $\alpha$ …

Soit un automorphisme $u$ de $X$ induisant l'identité sur $\partial X$. Il peut se relever en un automorphisme $\tilde u^{(S)}$ de $\widetilde{X}$ unique induisant l'identité sur $\widetilde{X}|S$ — alors son action sur $\widetilde{X}|S'$ est \struck{\ill{}} \add{induite} par une élément \struck{de} $n\in H_1(S')\simeq H_1(X)$.
\marginal{$S'=\partial X\setminus S$}
Si on \struck{change} \add{remplace} $S$ par $S'$, on trouve $\tilde u^{(S')}$, tel que
\[
\tilde u^{S'}=\tau_n^{-1}\tilde u^{S}=\tau_{-n}\tilde u^{S}
\]
\struck{\ill{}} quand on renverse : $S$ donne $\tau_{-n}$. Donc il y a une \uncertain{invariant}) qui vit dans
\[
H_1(X)\wedge_{\{\pm1\}}\varepsilon_0(X)=\bigl(\mathbb{Z}\wedge_{\pm1}\varepsilon_1(X)\bigr)\wedge_{\pm1}\varepsilon_0(X)
\simeq\mathbb{Z}\wedge_{\pm1}\bigl(\varepsilon_1(X)\wedge\varepsilon_0(X)\bigr)
\simeq\mathbb{Z}\wedge_{\pm1}\omega(X).
\]

\page{153}
\[
\text{\struck{$\tilde\rho^{3}(\tilde\tau'_\infty)=\tilde\omega_0^{-1}\tilde\tau'_\infty\tilde\omega_0=\tilde\omega_0^{-1}\tilde\tau'_\infty\tilde\omega_0\tilde\tau'^{-1}_\infty\tilde\tau'_0$}}
\]
\[
\text{\struck{$=\tilde\omega_0^{-2}\tilde\tau'_0=\ldots$}}
\]
\note{ligne biffée, en tête de page ; sa fin n'est pas sûre. Elle renvoie à des notations ($\tilde\rho$, $\tilde\omega_0$, $\tilde\tau'_\infty$) qui ne sont pas introduites dans ce lot.}
\[
\tau'',\ \tau,\quad \xi_0,\ \xi_1,\ \xi_2
\]
\[
\begin{array}{lll}
\tau'',\ \xi_0,\ \xi_1 & \text{engendrent} & \mathrm{Gl}(2,\mathbb{Z})\\
\tau'',\ \xi_2,\ \xi_1^{-1} & \text{———} & \mathrm{Gl}(2,\mathbb{Z})\\
\tau'',\ \xi_1 & \text{———} & D_{\xi_1}
\end{array}
\]
\note{devant « engendrent », un mot biffé. Le tiret long répète « engendrent ».}
\marginal{$\tau$ permute ces deux ss-groupes $\mathrm{Gl}(2,\mathbb{Z})$ et opère sur $D_{\xi_1}$}
\[
\mathcal{T}_{1,2}\ \overset{?}{\simeq}\ \mathrm{Gl}(2,\mathbb{Z})\underset{D_\infty}{*}\mathrm{Gl}(2,\mathbb{Z})
\]
\note{un « ! » est écrit au-dessus de l'indice de $\mathcal{T}$, un « ? » au-dessus du signe d'isomorphie.}

\drawing{un grand cercle (une sphère) ; trois arcs descendent depuis le haut, issus d'un même voisinage du pôle, et s'arrêtent vers le milieu.}
$3g$ cylindres qui sont deux à deux en position standard (deux cylindres polaires de \uncertain{même} pôle, ou un cylindre polaire et un transverse relatifs au même axe équatorial) ou ne se rencontrent pas.
\[
2g+2\,\frac{g(g-1)}{2}=2g\Bigl(1+\frac{g-1}{2}\Bigr)=2\,\frac{g(g+1)}{2}
\]
paires de cylindres en position standard$^{4}$, et les autres$^{5}$ \uncertain{paires} sont disjointes.
\note{les exposants 4 et 5 sont sur la page ; leur sens n'est pas indiqué. Une étiquette imprimée est collée sur ces deux lignes et en masque une partie ; « paires de cylindres » se lit en partie sous l'étiquette.}

\page{155}
\drawing{en haut, deux figures d'un tore (ou d'une sphère à anse) traversé par un cylindre horizontal. Sur la première sont portées les étiquettes $X'$ (deux fois), $\xi'$, $X_0^{*}$ et, sur le cylindre, $\xi''$ et une flèche marquée $5$ \uncertain{ou} $S$ ; la seconde montre le même cylindre traversant une figure ovale, vue de face.}

\[
\tau',\ \tau'',\ \tau^{*}\qquad \text{\struck{$\xi$}}\ \ \xi^{*}\ \ \bar\xi'
\]
\note{la lecture des exposants sur la seconde ligne est incertaine ; le premier symbole est surchargé.}

\drawing{une surface ovale découpée par des arcs (méridiens et un équateur), avec, sur une face, deux flèches parallèles de sens opposés et une flèche en pointillé. À droite : « 4 rectangles, 10 arêtes, 4 sommets ».}
\drawing{un petit cercle d'où partent trois segments ; entre deux d'entre eux, un secteur hachuré.}
\[
\tau'^2=\tau''^2=\tau^2=\text{\struck{\ill{}}}=1,\qquad (\tau''\tau)^2=(\tau\tau')^2=(\tau'\tau'')^2=1
\]
\note{le « $=1$ » de la première égalité est restitué d'après la suite de la ligne ; le passage biffé est illisible.}
\[
\begin{aligned}
&\tau(\xi_0)=\xi_2\qquad(\tau(\xi_2)=\tau(\xi_1))\\
&\tau(\xi_1)=\xi_1^{-1}\\
&\tau'(\xi_0)=\tau''(\xi_0)=\xi_0^{-1}\\
&(\Rightarrow\ \tau'(\xi_2)=\tau''(\xi_2)=\xi_2^{-1})\\
&\tau'(\xi_1)=\tau''(\xi_1)=\xi_1^{-1}\\
&[\xi_0,\xi_2]=1
\end{aligned}
\]
\note{la parenthèse « $(\tau(\xi_2)=\tau(\xi_1))$ » se lit ainsi sur la page ; on attendrait $\tau(\xi_2)=\xi_0$. Dans la ligne marquée $\Rightarrow$, l'indice du premier $\xi$ est surchargé ($1$ sur $2$).}
\[
\rho_0\overset{\text{déf}}{=}(\xi_1\xi_0)^{-1}=\xi_0^{-1}\xi_1^{-1},\qquad
\tau(\rho_0)=(\xi_1^{-1}\xi_2)^{-1}\overset{\text{déf}}{=}\rho_2=\xi_2^{-1}\xi_1
\]
\note{ces deux formules, à droite, sont séparées du reste par un trait courbe qui borde aussi la ligne horizontale au-dessous.}
\[
\xi_0\underbrace{\xi_1\xi_0}_{\rho_0^{-1}}=\underbrace{\xi_1\xi_0}_{\rho_0^{-1}}\xi_1
\ \Longrightarrow\ (\underbrace{\xi_1\xi_0}_{\rho_0^{-1}})^3=(\underbrace{\xi_0\xi_1\xi_0}_{\sigma_0})^2
\]
\[
(\xi_1\xi_0)^6=1,\qquad (\xi_1\xi_0)^6=(\xi_0\xi_1\xi_0)^4
\]
\note{les deux dernières relations sont entourées d'un ovale, d'où part une flèche vers la ligne suivante. Dans la première formule, le terme de gauche se lit $\xi_0\,\xi_1\xi_0$ ; on attendrait $\xi_0\xi_1\xi_0=\xi_1\xi_0\xi_1$, la relation de tresse.}
\[
\xi_2\xi_1^{-1}\xi_2=\xi_1^{-1}\xi_2\xi_1^{-1}
\]

\page{157}
\[
\begin{array}{lll}
\mathrm{Sl}(2,\mathbb{Z}) & \rho,\sigma & \bigm| \ \rho^3=\sigma^2,\ \sigma^4\,(=\rho^6)=1\\[1ex]
& \xi_0\,\xi_1 & \bigm| \ \xi_0\xi_1\xi_0=\xi_1\xi_0\xi_1,\ (\xi_1\xi_0)^6=1\\[2ex]
\mathrm{Gl}(2,\mathbb{Z}) & \rho,\sigma,r_0 & \bigm| \ \rho^3=\sigma^2,\ \sigma^4\,(=\rho^6)=1,\ r_0^2=1,\\
& & \phantom{\bigm|}\ r_0(\rho)=\rho^{-1},\ r_0(\sigma)=\sigma^{-1}\\[1ex]
& r_0,r_\infty,r_1 & \bigm| \ r_0^2=r_1^2=r_\infty^2=1,\ (r_0r_\infty)^2=(r_1r_0)^3,\ (r_0r_\infty)^4=1\\[1ex]
& r_0,\xi_0,\xi_1 & \bigm| \ r_0^2=1,\ r_0(\xi_0)=\xi_1^{-1},\ r_0(\xi_1)=\xi_0^{-1},\\
& & \phantom{\bigm|}\ \xi_0\xi_1\xi_0=\xi_1\xi_0\xi_1,\ (\xi_1\xi_0)^6=1
\end{array}
\]
\note{dans la présentation de $\mathrm{Sl}(2,\mathbb{Z})$ par $\xi_0,\xi_1$ comme dans la dernière de $\mathrm{Gl}(2,\mathbb{Z})$, une relation $(\xi_1\xi_0)^3(\xi_0\xi_1\xi_0)^{2}=1$ est biffée et remplacée par la relation de tresse. Dans la présentation par $r_0,r_\infty,r_1$, l'exposant de $(r_0r_\infty)$ est surchargé ($2$ sur $3$ \uncertain{ou} l'inverse), et les indices des générateurs sont surchargés eux aussi. Dans la dernière ligne, le second membre de la relation de tresse se lit $\xi_1\xi_2\xi_1$ ; on attend $\xi_1\xi_0\xi_1$.}

$\tilde\xi_0,\ \tilde\xi_1,\ \tilde r_0$ relèvements dus au fait que $\xi',\xi'',r_0$ fixent le pt base au-dessus de $\partial X$ ($\xi',\xi''$ sont $=$ l'identité sur $\partial X$)

On définit ensuite
\[
\begin{aligned}
\tilde\rho&=(\tilde\xi_1\tilde\xi_0)^{-1}=\tilde\xi_0^{-1}\tilde\xi_1^{-1}\quad\text{i.e.}\quad\tilde\xi_1\tilde\xi_0=\tilde\rho^{-1}\\
\tilde\sigma&=\tilde\xi_0\tilde\rho^{-1}\qquad\text{d'où}\quad\tilde\xi_0=\tilde\sigma\tilde\rho\\
&=\tilde\xi_0\tilde\xi_1\tilde\xi_0
\end{aligned}
\]
\note{à droite de la première ligne, un début de formule biffé ; sous la ligne $\tilde\xi_0=\tilde\sigma\tilde\rho$, une formule biffée et illisible.}

\struck{a-t-on $\tilde\rho(\tilde\xi_0)=\tilde\sigma(\tilde\xi_0)$ ?}

On a
\[
\begin{aligned}
&\tilde\sigma(\tilde\xi_0)=\sigma(\tilde\xi_0)=\tilde\xi_1\\
&\text{d'où}\quad \tilde\rho(\tilde\xi_0)=\tilde\sigma^{-1}(\tilde\xi_0)=\sigma^{-1}(\tilde\xi_0)=\sigma(\tilde\xi_0)=\tilde\xi_1
\end{aligned}
\]
\[
\tilde\rho^{-3}=(\tilde\xi_1\tilde\xi_0)^3,\qquad \tilde\sigma^2=(\tilde\xi_0\tilde\xi_1\tilde\xi_0)^2
\]
\marginal{on $\tilde\rho=\tilde\sigma^{-1}\tilde\xi_0$, d'où $\tilde\xi_1=\tilde\rho\tilde\sigma$ et $\tilde\sigma=\tilde\rho^{-1}\tilde\xi_1=\tilde\xi_1\tilde\xi_0\tilde\xi_1$}
\note{cette remarque, à droite, est isolée du reste par un trait.}
A-t-on
\[
\tilde\rho^{-3}=\tilde\sigma^2\quad\text{i.e.}\quad
\tilde\xi_1\tilde\xi_0\tilde\xi_1\tilde\xi_0\tilde\xi_1\tilde\xi_0=\tilde\xi_0\tilde\xi_1\tilde\xi_0^{2}\tilde\xi_1\tilde\xi_0\ \ ?
\]
\marginal{oui}
\note{sous « $\tilde\rho^{-3}=\tilde\sigma^2$ », une accolade porte la mention $\tilde\omega_0$.}
On voudrait de plus que
\[
\tilde\rho^{-6}=\tilde\sigma^4
\]
soit le générateur canonique du $\operatorname{Ker}(\mathcal{T}^{+}_{1,1}\to\mathcal{T}_{1,1})$
\note{la notation $\mathcal{T}^{+}_{1,1}$ est une lecture ; un signe surchargé précède $\mathcal{T}$.}

\page{159}
\[
\text{\struck{$\xi_0\ldots\tilde\xi_0$}}
\]
\[
\tilde\rho',\qquad \tilde\rho''=\sigma(\tilde\rho')
\]
\drawing{à gauche, un repère orthonormé d'axes $e'$ (horizontal) et $e''$ (vertical), avec un carré centré à l'origine, de sommets $a_0$ (en bas à gauche), $a_1$, $a_2$, $a_3$ ; les côtés $a_0a_1$ et $a_0a_3$ sont repassés en gras, la demi-diagonale de $a_0$ à l'origine est en pointillé, une flèche courbe marque une rotation d'un quart de tour dans le sens direct. L'axe horizontal porte la mention « $\tau'=\mathrm{c}^{\mathrm{te}}$ », l'axe vertical « $\tau''=\mathrm{c}^{\mathrm{te}}$ ».}
\note{« $\tau'=\mathrm{c}^{\mathrm{te}}$ » : lecture probable ; on comprend que l'axe est l'ensemble des points fixes de $\tau'$.}
\[
\begin{cases}
\tau'e'=e' & \tau''e'=-e'\\
\tau'e''=-e'' & \tau''e''=e''\\
\sigma e'=e'' & \sigma e''=-e'
\end{cases}
\]
\[
\tau'\tau''=\tau''\tau'=\sigma^2=\omega_0,\qquad \sigma^4=\omega_0^2=1
\qquad
\begin{cases}\omega_0e'=-e'\\ \omega_0e''=-e''\end{cases}
\]
\note{après $\sigma^2$, un signe surchargé ; sous ces lignes, deux lignes biffées et noircies, illisibles.}
\[
\sigma\tau'=\tau'_\infty=r_0\ :\quad e'\mapsto e'',\ \ e''\mapsto e'
\]
\[
\sigma\tau'=\tau'\sigma^{-1}=\sigma^{-1}\tau''=\tau''\sigma\qquad(\tau'=\sigma^2\tau'')
\]
\note{ces égalités sont écrites en colonne sous $\sigma\tau'$, reliées par des signes $=$ verticaux.}
\[
\begin{aligned}
\tau''&\longmapsto\tau_\infty=r_\infty\\
\sigma&\longmapsto\sigma=\sigma_\infty\\
\xi'&\longmapsto\xi_0\\
\xi''=\sigma(\xi')&\longmapsto\xi_1
\end{aligned}
\qquad\qquad
\begin{aligned}
\xi'(e')&=e'\\
\xi'e''&=e''-e'
\end{aligned}
\]
\marginal{$\xi'$ était $\rho'$ avant}
\note{dans $\xi'\mapsto\xi_0$ et dans les deux formules de droite, le $\xi'$ est écrit par-dessus un $e$ ou un $\rho$ ; la note marginale le confirme.}
\[
\xi_0=\sigma\rho,\quad \xi_1=\rho\sigma,\qquad \xi_1\xi_0=\rho\sigma^2\rho=-\rho^2=\rho^{-1}
\]
\note{la formule $\xi_1\xi_0=-\rho^2=\rho^{-1}$ est écrite ainsi ; l'étape intermédiaire $\rho\sigma^2\rho$ est restituée par le transcripteur pour la lisibilité, elle n'est pas sur la page. Au-dessous, une ligne biffée : $\xi_1\xi_0=\omega_0\xi_1\xi_0=\ldots$}
\[
\text{\struck{$\xi_0\ \xi_1\ \xi_0=\rho\sigma$}}\qquad
\begin{cases}
\sigma(\xi')=\xi''\\
\tau'(\xi')=\tau''(\xi')=\xi'^{-1}\\
\tau'(\xi'')=\tau''(\xi'')=\xi''^{-1}
\end{cases}
\]
\[
\begin{cases}
\rho=(\xi_1\xi_0)^{-1}=\xi_0^{-1}\xi_1^{-1}\\
\sigma_\infty=\xi_0\rho^{-1}=\xi_0\xi_1\xi_0=\rho^{-1}\xi_1=\xi_1\xi_0\xi_1\\
r_0(\xi_0)=\xi_1^{-1}\\
r_0(\xi_1)=\xi_0^{-1}
\end{cases}
\qquad
\begin{bmatrix}
r_0(\rho)=\xi_1\xi_0=\rho^{-1}\\
r_0(\sigma)=\xi_1^{-1}\xi_0^{-1}\xi_1^{-1}=\xi_0^{-1}\xi_1^{-1}\xi_0^{-1}
\end{bmatrix}
\]
\note{devant $r_0(\xi_0)$, un mot biffé ; au-dessus de $\rho^{-1}\xi_1$, une formule biffée.}
\[
\text{i.e.}\quad \xi_1\xi_0\xi_1=\xi_0\xi_1\xi_0
\]
\[
\text{\struck{$\rho\sigma\sigma\rho\rho\sigma=\sigma\rho\rho\sigma\sigma\rho$}}
\]
\[
\rho\sigma^2\rho^2\sigma=\sigma\rho^2\sigma^2\rho,\qquad -\rho^3\sigma=-\sigma\rho^3
\]
\note{sous la ligne biffée, une seconde ligne biffée en $\xi_0,\xi_1$, illisible. Sous $-\rho^3\sigma=-\sigma\rho^3$, les deux membres sont soulignés d'accolades portant chacune $\sigma$ (au sens : $\rho^3=\sigma^2$).}

$\xi_0,\ \xi_1,\ r_0$
\note{fin du lot : la page suivante (p. 160) est un document administratif sans rapport. Le calcul sur les générateurs $\xi_0,\xi_1,r_0$ de $\mathrm{Gl}(2,\mathbb{Z})$ se poursuit peut-être au-delà du lot.}

\end{document}
